% Part: computability
% Chapter: recursive-functions
% Section: primes

\documentclass[../../../include/open-logic-section]{subfiles}

\begin{document}

\olfileid{cmp}{rec}{pri}
\olsection{Wilangan Prima}

Kuantifikasi mawa wates lan minimisasi mawa wates menehi
cukup akeh piranti kanggo nuduhake yen fungsi lan relasi
ing wilangan asli iku rekursif primitif. Umpamane, gatekna
relasi “$x$ mbagi $y$,” kang ditulis $x \mid y$. Relasi
$x \mid y$ bener yen $y$ dipara~$x$ tanpa
turah, yaiku yen $y$ iku asil ping-pingan
sawenehing wilangan wutuh karo~$x$. (Yen relasi
iki ora bener, mula kanggo pembagi kang positif,
turah nalika $y$ dipara $x$
iku $> 0$; kita nulis $x \nmid y$.)
% OLPL-125: Sumber mbalikke urutan kang dibagi lan pembagi ing
% katerangan turah; turah sing trep yaiku saka y dipara x.
Kanthi tembung liya, $x \mid y$ yen lan mung yen ana~$z$
kanthi $x \cdot z = y$. Kita bisa milih saksi~$z$
kang $\leq y$ yen saksi ana. Dadi $x \mid y$
yen lan mung yen ana $z \le y$ kanthi $x \cdot z = y$.
% OLPL-126: Pratelan sumber "saben saksi z <= y" salah yen
% x=y=0; saksi z=0 tetep ana lan nyukupi wates.
Kita bisa netepake relasi $x \mid y$ kanthi kuantifikasi
eksistensial mawa wates saka $=$ lan perkalian:
\[
x \mid y \defiff \bexists{z \leq y}{(x \cdot z) = y}.
\]
Mula kita wis nuduhake yen $x \mid y$ rekursif primitif.

Wilangan asli~$x$ diarani \emph{prima} yen dudu $0$
utawa $1$ lan mung bisa dibagi dening $1$ lan
awake dhewe. Kanthi tembung liya, kanggo wilangan
prima, yen $y \mid x$, mesthi $y = 1$ utawa~$y=x$.
Kanggo mriksa apa $x$~prima, kita mung perlu mriksa
$y \mid x$ kanggo kabeh $y \le x$, amarga yen
$y > x$, mesthi~$y \nmid x$. Dadi relasi
$\fn{Prime}(x)$, kang bener yen lan mung yen $x$
prima, bisa ditegesi kanthi
\[
\fn{Prime}(x) \defiff x \geq 2 \land \bforall{y \leq x}{(y \mid x \lif y
  = 1 \lor y = x)}
\]
lan mula rekursif primitif.

Wilangan prima yaiku $2$, $3$, $5$, $7$, $11$,
lan sateruse. Gatekna fungsi $p(x)$ kang ngasilake
prima kaping $x$ ing rerangken kasebut, yaiku
$p(0) = 2$, $p(1) = 3$, $p(2) = 5$, lan sateruse.
(Kanggo nggampangake, kita kerep nulis $p(x)$
minangka $p_x$ ($p_0=2$, $p_1=3$, lan sateruse.)

Yen kita nduweni fungsi
$\fn{nextPrime(x)}$, kang ngasilake wilangan prima
pisanan kang luwih gedhe tinimbang~$x$, $p$~bisa
gampang ditegesi nganggo rekursi primitif:
\begin{align*}
  p(0) & = 2\\
  p(x+1) & = \fn{nextPrime}(p(x))
\end{align*}
Amarga $\fn{nextPrime}(x)$ iku $y$ paling cilik
kang nyukupi $y > x$ lan $y$~prima, fungsi iki
bisa gampang diitung kanthi panggolekan tanpa wates.
Nanging fungsi iki uga bisa ditegesi kanthi minimisasi
mawa wates, amarga asil saka Euclid: mesthi ana
wilangan prima ing antarane $x$ lan $\fact{x}+1$.
\[
  \fn{nextPrime(x)} =
  \bmin{y \leq \fact{x}+1}{(y > x \land \fn{Prime}(y))}.
\]
Iki nuduhake yen $\fn{nextPrime}(x)$ lan mula $p(x)$
(ora mung komputabel nanging) rekursif primitif.

(Yen kepengin ngerti, iki bukti ringkes kanggo
teorema Euclid. Kanggo wilangan asli luwih gedhe
tinimbang siji, upamana $p_n$ prima paling gedhe
kang $\le x$, lan gatekna asil kali $p =
p_0\cdot p_1 \cdot \dots \cdot p_n$ saka kabeh
prima~$\le x$. Yen wilangan asline nol utawa siji,
prima loro langsung menehi kesimpulan kang padha.
% OLPL-127: Bukti nganggo "prima paling gedhe <= x" mung ana
% yen x paling ora loro; kasus nol lan siji dipriksa langsung.
$p+1$ iku prima utawa ana prima ing antarane
$x$ lan~$p+1$. Ngapa? Upamana $p+1$ dudu prima.
Banjur ana prima $q \mid p+1$ kang $q < p+1$.
Ora ana prima $\le x$ kang mbagi $p+1$. (Miturut
definisi~$p$, saben prima $p_i \le x$ mbagi~$p$,
yaiku kanthi turah~$0$. Mula saben prima $p_i \le x$
mbagi $p+1$ kanthi turah~$1$, dadi $p_i \nmid p+1$.)
Mula $q$ iku prima $> x$ lan $< p+1$. Lan
$p \le \fact{x}$, dadi ana prima $> x$ lan
$\le \fact{x}+1$.)

\begin{prob}
Netepna pambagian wilangan wutuh $d(x, y)$ nganggo
minimisasi mawa wates.
\end{prob}

\end{document}
